
\begin{theorem}[Support of the original product site map]
    \label{thm:peps_original_region_projector_absorption}
    \lean{TNLean.PEPS.regionPhysicalProductMatrix_mul_regularLocalProjector}
    \leanok
    For original regular $G$-injective site tensors on a finite region, let
    $A$ be their product physical map and $P$ the product of their local
    averaging projectors. Then $AP=A$.
    This is the local-invariance identity used in
    \cite{Schuch2010PEPS}, source lines 1765--1820.
\end{theorem}
\begin{proof}\leanok
    Each original site map absorbs its local averaging projector.
    Multiplication of the product maps gives the regional identity.
\end{proof}

\begin{theorem}[Return measurement on the original physical block]
    \label{thm:peps_original_charge_pair_return_measurement}
    \lean{TNLean.PEPS.exists_regularPhysicalChargePairReturnMeasurement}
    \leanok
    \uses{thm:peps_reference_return_actual_column,
      thm:peps_original_region_projector_absorption,
      thm:peps_scaled_projection_transport}
    Let $R$ be a finite region with a chosen spanning tree and root and two
    distinct internal bonds $e_0,e_1$. Assume that every original site map
    is regularly $G$-isometric. Fix a unitary irreducible character $\chi$
    of dimension $d$ and $p\in G$.
    There is a complete positive binary projection measurement $(Q_0,Q_1)$
    on the original spins, chosen before $k\in G$ and every actual boundary
    label $\theta$, with the following action. Let $\psi_{k,\theta}$ be the
    literal original open contraction with weight
    $\chi(p\eta_{e_1}^{-1}k^{-1}\eta_{e_0})$, and let
    $\psi_{1,\theta}$ be the initial correlated column. Then
    \begin{align}
        Q_1\psi_{k,\theta}&=\frac{\chi(k^{-1})}{d}\psi_{1,\theta},\notag\\
        Q_0\psi_{k,\theta}&=\psi_{k,\theta}
          -\frac{\chi(k^{-1})}{d}\psi_{1,\theta}.\notag
    \end{align}
    No decomposition of a global physical state is assumed. This is an
    auxiliary finite-region return operation for
    \cite[Section~6.6]{Schuch2010PEPS}, source lines 2582--2615.
\end{theorem}
\begin{proof}\leanok
    The local $G$-isometries give $A^\dagger A=cP$ with $c>0$, while
    regular invariance gives $AP=A$. The actual extended reference return
    projector is Hermitian and idempotent and commutes with $P$.
    Its transport $c^{-1}ADA^\dagger$ and the physical complement are a
    complete binary measurement. Apply its intertwining identity to the
    already derived canonical correlated column, then recover the literal
    original weighted contraction.
\end{proof}

\begin{theorem}[Return measurement on six translated torus spins]
    \label{thm:peps_torus_six_spin_return_measurement}
    \lean{TNLean.PEPS.IsGIsometric.exists_torusPhysicalChargePairReturnMeasurement}
    \leanok
    \uses{thm:peps_original_charge_pair_return_measurement,
      thm:peps_torus_physical_charge_pair_creation}
    On the actual translated two-column, three-row torus patch, the preceding
    measurement acts on the six original spins. The reference bonds, spanning
    tree and root are derived from the native geometry. The same measurement
    precedes every $k$ and $\theta$, including patches crossing either seam.
    The horizontal period is at least three and the vertical period at least
    four. This proves the local return operation on the literal correlated
    bond sum; identifying a complete prescribed geometric braid with that
    input remains separate. See \cite{gap:rmp_peps_quantum_double_g_isometry}.
\end{theorem}
\begin{proof}\leanok
    Apply the finite-region theorem to the actual six-site path tree and its
    two distinct horizontal chords. Native coordinate sorting retains the
    original boundary labels through both seams.
\end{proof}
